% =====================================================================
%  REMERCIEMENTS
% =====================================================================
\chapter*{Remerciements}
\addcontentsline{toc}{chapter}{Remerciements}
\markboth{Remerciements}{Remerciements}

Je souhaite exprimer ma profonde reconnaissance à toutes les personnes qui ont,
de près ou de loin, contribué à la réussite de mon stage de fin d'études et à
l'aboutissement de ce travail.

Je remercie tout d'abord mes parents et ma famille, dont le soutien constant, la
confiance et les encouragements indéfectibles m'ont permis d'avancer avec
sérénité tout au long de mes études et de ce stage. Leur présence à mes côtés a
été une source inestimable de force et de motivation, y compris dans les moments
de doute.

J'exprime ensuite mes sincères remerciements à mon tuteur universitaire,
M.~Taoufik M'hammedi, pour son encadrement attentif, ses conseils avisés et la
disponibilité dont il a fait preuve tout au long de cette expérience. Je tiens
également à remercier M.~Nourdin Yaakoubi pour l'honneur qu'il me fait en
évaluant ce travail en tant que membre du jury, ainsi que pour l'attention portée
à mon parcours.

Mes remerciements vont également à mes tuteurs en entreprise, M.~Alexis Ster,
team lead, et M.~Maxime Villermin, chef de projet, pour la confiance qu'ils
m'ont accordée dès les premières semaines, la richesse de leurs enseignements et
la qualité de leur accompagnement. Leur rigueur professionnelle, alliée à leur
bienveillance, a grandement contribué à la réussite de mon intégration et à
l'enrichissement de mes compétences. Je remercie tout particulièrement Nathan
Richez, référent technique, dont les points hebdomadaires ont guidé mes choix de
conception et m'ont permis de progresser rapidement sur des technologies que je
découvrais.

Je souhaite aussi adresser mes vifs remerciements à l'ensemble de l'équipe Group
Collaborative Solutions et, plus largement, aux collaborateurs de la direction
Custom Applications de VESI, pour leur accueil chaleureux, leur disponibilité et
la qualité des échanges quotidiens. Leur expertise et leur esprit d'équipe ont
été des sources précieuses d'apprentissage et d'inspiration.

Je remercie également les enseignants de l'ENSIM dont les cours ont directement
nourri le travail présenté dans ce rapport, en particulier en bases de données,
en génie logiciel et en développement web. C'est la solidité de ces fondations
qui m'a permis de monter en compétence sans blocage majeur sur des technologies
que je n'avais pas pratiquées auparavant.

Enfin, je remercie chaleureusement toutes les personnes avec lesquelles j'ai eu
l'occasion d'échanger ou de collaborer durant ces six mois. Chaque interaction a
représenté une opportunité d'apprendre et de progresser, rendant cette
expérience particulièrement formatrice et marquante.
